% Generated by tools/edn_latex.py from reports/edn.md.
% Do not edit: change reports/edn.md and run `make edn`.
\ednentry{
  id = {EDN-50},
  title = {One-hot encoding for the Ridge baseline only, against EDN-02's ``no one-hot''},
  date = {2026-09-30},
  milestone = {M3: Quality Assurance},
  topic = {Feature Engineering, Modelling Approach},
  participants = {@lukas2510},
  decision = {The \texttt{b1} Ridge variant one-hot encodes its categoricals, with \texttt{handle\_\allowbreak{}unknown=\allowbreak{}\textquotedbl{}infrequent\_\allowbreak{}if\_\allowbreak{}exist\textquotedbl{}} and \texttt{min\_\allowbreak{}category\_\allowbreak{}rows: 5} as the threshold below which levels share one column. This is confined to \texttt{b1}: the LightGBM variants and the CatBoost challenger take the categoricals natively, as \texttt{category} codes over the levels the \texttt{features} stage fixed, and no one-hot column exists anywhere near them.},
  rationale = {\emph{Alternatives considered:}
\begin{itemize}[leftmargin=*, topsep=2pt]
\item \textbf{Option A (chosen): one-hot, confined to B1.}
\newline Pros: it is the standard encoding for a linear model, it needs no target information, and it keeps B1 able to do the job it exists for -- telling a Porsche from a Dacia as part of an interpretable depreciation baseline. \texttt{handle\_\allowbreak{}unknown=\allowbreak{}\textquotedbl{}infrequent\_\allowbreak{}if\_\allowbreak{}exist\textquotedbl{}} is also what makes an unseen level an answer rather than an error, which EDN-18 requires of the serving path.
\newline Cons: it is the encoding EDN-02 rules out, so a reader who takes that rule as project-wide sees a violation. The column count grows with the level count: 503 columns at real scale before any folding, against 16 input features.
\item \textbf{Option B: target encoding with cross-fitting.}
\newline Pros: one column per categorical however many levels it has, and it usually beats one-hot for a linear model on high-cardinality data.
\newline Cons: it encodes the target into a feature, so it leaks unless it is cross-fitted, and cross-fitting means a second resampling scheme inside the stage, a second thing to persist and a second thing that can differ between training and serving. For a baseline whose purpose is to be simple and auditable that is the wrong trade, and the leak is the kind of defect this project would only find by not finding it.
\item \textbf{Option C: drop the categoricals from B1 and fit on the numerics alone.}
\newline Pros: no encoding decision at all, and the ladder still has a linear baseline.
\newline Cons: it destroys the baseline. A depreciation model that cannot see the make is not a weaker version of B1, it is a different and much worse thing, and SC-03's comparison against B0 would be against a baseline nobody would have built.
\item \textbf{Option D: read EDN-02 as project-wide and drop B1 from the ladder.}
\newline Pros: no apparent contradiction to explain.
\newline Cons: problem-spec section 7 names B1 as a baseline, and the interpretable linear reference is what makes the tree models' gain legible. Removing a baseline to preserve the letter of a rule about a different model family is the wrong direction.
\end{itemize}
\par
\emph{Why:} EDN-02 is a decision about the \emph{model family}, and ``categoricals stay categorical, no one-hot'' is stated there as one of the reasons gradient boosting was chosen over the alternatives -- LightGBM and CatBoost handle them natively, so the pipeline does not have to blow up the matrix. It is not a constraint on how a linear baseline encodes its inputs, because a linear model has no native handling to use: the choice is one-hot, target encoding or nothing. Recording this keeps the next reader from either ``fixing'' B1 to match the rule or reading the rule as broken.
\par
Two details of the encoding are decisions in their own right and are measured rather than assumed.
\par
Absence is a level. Categoricals are mapped to a literal \texttt{\textquotedbl{}\_\allowbreak{}\_\allowbreak{}missing\_\allowbreak{}\_\allowbreak{}\textquotedbl{}} value before encoding, deterministically at fit and at predict time, so an absent value gets a coefficient like any other level instead of failing the fit or being dropped. That is the same treatment EDN-15 asks for, expressed in the only way a linear model can express it.
\par
\texttt{min\_\allowbreak{}category\_\allowbreak{}rows} is an absolute row count rather than a share, so that what the threshold means does not change as the training set grows. Measured on the real snapshot (60,378 training rows, MdAPE on the 19,665 test rows): no folding gives 503 encoded columns at 9.48\,\%, 5 gives 404 at 9.52\,\%, 30 gives 294 at 9.71\,\%, 100 gives 187 at 10.61\,\% and 1,000 gives 68 at 13.45\,\%. So folding costs accuracy monotonically, and the reason to fold is not accuracy.
\par
What the reason is, read off the fitted encoder rather than assumed: at 5, folding touches \textbf{2 of the 8 categoricals}. \texttt{model} folds 100 of its 438 levels into one infrequent column and \texttt{fuel\_\allowbreak{}category} folds 1. The other six have no infrequent group at all, and \texttt{make} and \texttt{body\_\allowbreak{}type} have no \texttt{\_\allowbreak{}\_\allowbreak{}missing\_\allowbreak{}\_\allowbreak{}} level either, because the real training rows never leave those two empty. An unseen or an absent value in one of those six therefore encodes as all zeros rather than into a group.
\par
EDN-18 is satisfied either way, and checked rather than assumed: masking each of the eight in turn, and giving each an unseen value in turn, the real \texttt{b1} answers with a finite positive price all sixteen times. It is simply satisfied by the all-zero encoding for six of the eight -- a coherent encoding of ``no level applies'' -- and not by the infrequent group the threshold produces. The threshold is what stops \texttt{handle\_\allowbreak{}unknown=\allowbreak{}\textquotedbl{}infrequent\_\allowbreak{}if\_\allowbreak{}exist\textquotedbl{}} from being pointless for \texttt{model}, which is the one column where rare levels are the norm, and nothing more.
\par
So the justification for 5 is what is left: folding is nearly free there, 0.04 pp, and it stops a level seen in a handful of rows from getting a coefficient fitted on those rows. At real scale that is 47 \texttt{model} levels seen exactly once, among the 100 the threshold folds. An infrequent group does not need 5 to exist -- 2 already gives one, folding 48 levels for 0.02 pp -- so 5 is a round number inside a flat region rather than a boundary, and the value is a parameter precisely so the ladder can revisit it.
\par
One claim that did \textbf{not} survive measurement and is recorded so nobody repeats it: the design this work followed justified a much larger threshold with a small-data pathology, a mid-range car with an unseen model and country predicted at 94,323 EUR. That did not reproduce here at either scale. On the fixture the same probe gives 13,993 EUR unfolded and 15,413 EUR at a threshold of 30; at real scale, 32,332 EUR and 33,104 EUR. The threshold is justified by the infrequent group and by the one-row-coefficient argument, not by a pathology this project observed.},
  ai = {Information seeking, Alternative generation, Alternative assessment, Recommendation},
  response = {Accepted with modifications},
  assessment = {AI spotted the apparent conflict with EDN-02 before writing the estimator and flagged it as an EDN candidate rather than resolving it silently, and its reading -- that EDN-02 is about the tree family -- is the one the entry records. It was overruled on the threshold: it recommended 30 on the strength of a measurement from another run and asserted the value was free at real scale. Sweeping it here showed folding is not free and that its pathology did not reproduce, so the value is 5 and the entry says what the evidence actually supports.},
  aievidence = {Claude Code session on 2026-09-30 implementing issue \#37: the threshold sweep was run on the fixture and then on the real snapshot, and the ``unseen model and country'' probe was re-measured at both scales before the recommended value was changed.},
  evidence = {\href{https://github.com/mlops-2627q1-mds-upc/MLOPS_Recommenditos/blob/main/recommenditos/modeling/model.py}{\texttt{recommenditos/\allowbreak{}modeling/\allowbreak{}model.\allowbreak{}py}}, \texttt{RidgeModel}; the \texttt{min\_\allowbreak{}category\_\allowbreak{}rows} comment in \href{https://github.com/mlops-2627q1-mds-upc/MLOPS_Recommenditos/blob/main/params.yaml}{\texttt{params.\allowbreak{}yaml}}; \href{https://github.com/mlops-2627q1-mds-upc/MLOPS_Recommenditos/blob/main/docs/docs/model-card.md}{model card}, Training Procedure, Training; \texttt{tests/\allowbreak{}test\_\allowbreak{}model.\allowbreak{}py::test\_\allowbreak{}an\_\allowbreak{}unseen\_\allowbreak{}category\_\allowbreak{}is\_\allowbreak{}treated\_\allowbreak{}as\_\allowbreak{}missing\_\allowbreak{}not\_\allowbreak{}an\_\allowbreak{}error}; \hyperref[EDN-02]{EDN-02}; \hyperref[EDN-18]{EDN-18}; \href{https://github.com/mlops-2627q1-mds-upc/MLOPS_Recommenditos/issues/37}{issue \#37}.},
}
